% \section{Stationary Iterative Methods for Solving Linear Systems}

% General Iteration Scheme and Matrix Splitting

% For a linear system $A\mathbf{x}=\mathbf{b}$ with $A \in \mathbb{R}^{n\times n}$, stationary iterative methods generate a sequence of approximations ${\mathbf{x}^{(k)}}$ to the solution by a fixed linear recurrence. The idea is to split the coefficient matrix $A$ into two parts, $A = M - N$, where $M$ is chosen to be easily invertible (e.g. diagonal or triangular) and $N = M - A$  . Using this splitting, the system $A\mathbf{x}=\mathbf{b}$ can be rewritten as $M\mathbf{x} = N\mathbf{x} + \mathbf{b}$, leading to the iterative update formula (assuming $M$ is nonsingular):

% \begin{equation}
% \mathbf{x}^{(k+1)} = M^{-1}\Big( N\,\mathbf{x}^{(k)} + \mathbf{b}\Big)\,. \tag{1}
% \end{equation}

% This is the general stationary iteration scheme  . We start from an initial guess $\mathbf{x}^{(0)}$ and apply (1) repeatedly. The method is called stationary because the iteration matrix $B = M^{-1}N$ and the vector $M^{-1}\mathbf{b}$ remain the same at each step (in contrast to more advanced nonstationary methods where the update changes with $k$). The error at step $k$ is $\boldsymbol{\epsilon}^{(k)} = \mathbf{x}^{(k)} - \mathbf{x}$ (with $\mathbf{x}$ the true solution). Subtracting the true solution from (1) gives $\boldsymbol{\epsilon}^{(k+1)} = M^{-1} N,\boldsymbol{\epsilon}^{(k)} = B,\boldsymbol{\epsilon}^{(k)}$. Thus the error evolves as $\boldsymbol{\epsilon}^{(k)} = B^k ,\boldsymbol{\epsilon}^{(0)}$. For convergence from any initial guess, we require $\boldsymbol{\epsilon}^{(k)} \to \mathbf{0}$ as $k\to\infty$, which is equivalent to $B^k \to 0$. Necessary and sufficient condition: the spectral radius of $B$ must be less than 1 :

% \begin{equation}
    
% \rho(B) < 1 \implies B^k \to 0 \text{ as } k\to\infty\,. \tag{2}
% \end{equation}

% In fact, $\rho(B)<1$ is the condition for convergence of the iteration for any $\mathbf{x}^{(0)}$ . In practice, it is often difficult to compute $\rho(B)$, so one uses sufficient conditions: if one can find a matrix norm $|\cdot|$ such that $|B| < 1$, then $|B^k| \le |B|^k \to 0$ as $k\to\infty$, guaranteeing convergence (by the Banach fixed-point theorem). Classic criteria like diagonal dominance ensure this. For example, if $A$ is strictly diagonally dominant by rows, one can show $|B|_\infty < 1$ for the iteration matrix $B$ of both Jacobi and Gauss–Seidel methods, hence convergence . We will discuss specific convergence conditions for each method later on.

% The efficiency of a stationary method depends on how $A$ is split. The matrix $M$ should be easy to invert (solving $M\mathbf{y}=\mathbf{c}$ should be much cheaper than solving $A\mathbf{y}=\mathbf{c}$ directly) . Common choices are:
% 	•	$M =$ diagonal part of $A$ (leads to the Jacobi method),
% 	•	$M =$ triangular part of $A$ (leads to Gauss–Seidel),
% 	•	or more generally $M =$ diagonal part scaled by a factor plus one triangular part (leads to Successive Over-Relaxation, SOR).

% These choices yield the classical stationary methods discussed below. Each method results in a specific iteration matrix $B = M^{-1}N$. The rate of convergence is determined by $\rho(B)$ (smaller is better), and the method converges if $\rho(B)<1$. We now derive the Jacobi, Gauss–Seidel and SOR methods via different matrix splittings and examine their properties.

% Jacobi Method

% The Jacobi method is the simplest stationary iterative solver for $A\mathbf{x}=\mathbf{b}$. It splits $A$ into its diagonal $D$ and off-diagonal parts $L$ and $U$ (strictly lower and strictly upper triangular, respectively):

% \begin{equation}
    
% A = D + L + U,
% \end{equation}

% where $D = \operatorname{diag}(a_{11},a_{22},\dots,a_{nn})$. In the Jacobi splitting, we take $M = D$ and $N = -(L+U)$, so that $M - N = D - (-(L+U)) = D + L + U = A$. Plugging into (1), we get the Jacobi iteration:

% \begin{equation}
    
% \mathbf{x}^{(k+1)} = D^{-1}\Big(\mathbf{b} - (L+U)\,\mathbf{x}^{(k)}\Big)\,. \tag{Jacobi}
% \end{equation}

% This means each component $x_i$ is updated using only values from the previous iteration:

% \begin{equation}
    
% x_i^{(k+1)} = \frac{1}{a_{ii}}\Big(b_i - \sum_{j\neq i} a_{ij}\,x_j^{(k)}\Big), \qquad i=1,2,\ldots,n .
% \end{equation}

% In other words, one Jacobi iteration computes the new $x_i^{(k+1)}$ by “solving” the $i$-th equation of $A\mathbf{x}=\mathbf{b}$ for $x_i$, assuming all other components $x_j$ (with $j\ne i$) are fixed to their old values from iteration $k$. The method can be implemented in a vectorized way using the formula above, or component-wise as given. Because each $x_i^{(k+1)}$ uses only the old iterate $\mathbf{x}^{(k)}$, Jacobi updates can be done in parallel for all $i$ – one does not overwrite $\mathbf{x}^{(k)}$ until the entire new vector $\mathbf{x}^{(k+1)}$ is computed . This requires storing two vectors (for old and new values), but allows straightforward parallelization across $n$ equations.

% Convergence and properties: The Jacobi iteration matrix is $B_J = D^{-1}(L+U) = I - D^{-1}A$. Its convergence requires $\rho(B_J)<1$ . A sufficient condition is that $A$ be strictly diagonally dominant (or irreducibly diagonally dominant) , i.e. $|a_{ii}| > \sum_{j\neq i}|a_{ij}|$ for each $i$. In that case $|B_J|_\infty < 1$ and Jacobi converges. In fact, diagonal dominance is often used as an easy check for convergence of both Jacobi and Gauss–Seidel  . Jacobi can sometimes converge even if these conditions are not met , but convergence is not guaranteed in general. Notably, Jacobi does not always converge for symmetric positive definite (SPD) matrices : SPD is not a sufficient condition by itself. For example, the SPD matrix

% \begin{equation}
    
% A=\begin{pmatrix}29 & 2 & 1\\ 2 & 6 & 1\\ 1 & 1 & \frac{1}{5}\end{pmatrix}
% \end{equation}

% has a Jacobi iteration matrix with $\rho(B_J)\approx 1.0661>1$, so the Jacobi method diverges in this case . (In technical terms, for Jacobi to converge on an SPD matrix $A$, $A$ must be diagonally dominant enough; one sufficient condition is that both $A$ and $2D - A$ are SPD .) In practice, Jacobi is usually slower to converge than Gauss–Seidel for the same matrix if both converge. However, Jacobi’s simplicity and parallelism can be advantages. It is sometimes used as a smoother in multigrid methods or as a preconditioner, despite its relatively slow convergence, because each iteration is cheap and easily parallelizable.

% Gauss–Seidel Method

% The Gauss–Seidel method improves upon Jacobi by using the most recently updated values as soon as they are available, instead of waiting for the next iteration. It corresponds to the splitting $A = (D+L) - U$, i.e. $M = D+L$ (the lower triangular part of $A$ including the diagonal) and $N = U$ (the strictly upper part). The iteration (1) then becomes:

% \begin{equation}
    
% (D+L)\,\mathbf{x}^{(k+1)} = \mathbf{b} - U\,\mathbf{x}^{(k)}\, ,
% \end{equation}

% so that we solve a triangular system each step:

% \begin{equation}
    
% \mathbf{x}^{(k+1)} = (D+L)^{-1}\Big(\mathbf{b} - U\,\mathbf{x}^{(k)}\Big)\,. \tag{Gauss–Seidel}
% \end{equation}

% This can be written component-wise as a forward sweep through $i=1$ to $n$:

% \begin{equation}
    
% x_i^{(k+1)} = \frac{1}{a_{ii}}\Big(b_i \;-\; \sum_{j< i} a_{ij}\,x_j^{(k+1)} \;-\; \sum_{j> i} a_{ij}\,x_j^{(k)}\Big), \qquad i=1,2,\ldots,n  .
% \end{equation}

% Here we use the newest available values: when computing $x_i^{(k+1)}$, all “previous” components $x_1^{(k+1)},\dots,x_{i-1}^{(k+1)}$ have already been updated in the current iteration, while $x_{i+1}^{(k)},\dots,x_n^{(k)}$ are still from the old iteration. In matrix form, $L\mathbf{x}^{(k+1)} = \mathbf{b} - U\mathbf{x}^{(k)}$ (with $L=D+L$ including the diagonal) . This asynchronous update tends to incorporate new information faster than Jacobi. Practically, Gauss–Seidel can be done in-place (overwriting) a single vector as we iterate $i=1$ to $n$, since once $x_i^{(k+1)}$ is computed, we will not need $x_i^{(k)}$ again.

% Convergence and properties: The Gauss–Seidel iteration matrix is $B_{GS} = (D+L)^{-1}U$. Gauss–Seidel often converges under similar conditions to Jacobi, and in fact it converges under somewhat weaker conditions. A sufficient condition for convergence is again strict or irreducible diagonal dominance of $A$ . More importantly, if $A$ is symmetric positive definite, Gauss–Seidel is guaranteed to converge  (in contrast to Jacobi). In general, whenever the Jacobi method converges, the Gauss–Seidel method will also converge and usually at a faster rate  . The reason is that $\rho(B_{GS})\le \rho(B_J)$ for the same matrix, often strictly smaller . Intuitively, by using fresher data, Gauss–Seidel reduces the error more per iteration. In fact, for many SPD matrices one can show $\rho(B_{GS}) = [\rho(B_J)]^2$ (see Convergence Analysis below), suggesting Gauss–Seidel squares the convergence factor of Jacobi, which is a significant improvement.

% On the downside, Gauss–Seidel’s updates are sequential and depend on the ordering of the unknowns. You cannot update all $i$ in parallel, which makes it less suitable for parallel computing compared to Jacobi . (There are techniques like coloring or domain decomposition to parallelize Gauss–Seidel by grouping independent updates, but the method is inherently sequential in a given ordering.) Another consideration is that different orderings of the equations or variables can affect the convergence speed of Gauss–Seidel. For certain matrices, a reverse order (sweeping $i=n$ down to $1$, sometimes called backward Gauss–Seidel) might perform similarly; doing one forward and one backward sweep gives the Symmetric Gauss–Seidel iteration, often used in analysis and as a preconditioner. In any case, the simplicity of Gauss–Seidel (just a triangular solve each iteration) and its typically better convergence per step than Jacobi make it a popular choice for many problems, especially if one can afford sequential updates.

% Successive Over-Relaxation (SOR)

% Successive Over-Relaxation (SOR) is a one-parameter generalization of Gauss–Seidel that can accelerate convergence by over-correcting the solution at each step. The method introduces a relaxation factor $\omega>0$. For $\omega=1$, SOR reduces to the Gauss–Seidel method; for $\omega \neq 1$, the update is a weighted combination of the old iterate and the would-be Gauss–Seidel update. Formally, the SOR iteration can be derived by splitting $A$ as:
% \begin{equation}
    
% A = \frac{1}{\omega}D + L \;-\; \Big(\frac{1-\omega}{\omega}D + U\Big).
% \end{equation}

% Here $M = \frac{1}{\omega}D + L$ (a lower triangular matrix) and $N = \frac{1-\omega}{\omega}D + U$. Plugging into the general iteration (1) yields:

% \begin{equation}
    
% (D + \omega L)\,\mathbf{x}^{(k+1)} \;=\; \omega\,\mathbf{b} - [\,\omega U + (\omega - 1)D\,]\,\mathbf{x}^{(k)} ,
% \end{equation}

% so that

% \begin{equation}
    
% \mathbf{x}^{(k+1)} \;=\; (D+\omega L)^{-1}\Big(\omega\,\mathbf{b} - [\,\omega U + (\omega - 1)D\,]\,\mathbf{x}^{(k)}\Big)\,. \tag{SOR}
% \end{equation}

% This can be understood as solving a modified triangular system at each step. It is instructive to look at the element-wise formula for SOR (with the forward sweep assumption as in Gauss–Seidel):


% \begin{equation}
% x_{i}^{(k+1)} \;=\; (1-\omega)\,x_{i}^{(k)} \;+\; \frac{\omega}{a_{ii}}\Big(b_{i} - \sum_{j<i}a_{ij}\,x_{j}^{(k+1)} - \sum_{j>i}a_{ij}\,x_{j}^{(k)}\Big), \qquad i=1,2,\ldots,n .
% \end{equation}
% The term in parentheses is precisely the Gauss–Seidel update for $x_i$ (the value it would have if we did a normal Gauss–Seidel iteration), and SOR takes a weighted average between the old value $x_i^{(k)}$ and this new Gauss–Seidel value. If $\omega=1$, we take the new value fully (recovering Gauss–Seidel). If $1 < \omega < 2$, we over-relax by taking the new value and then extrapolating a bit further in that direction – hopefully converging faster. If $0 < \omega < 1$, this is under-relaxation, meaning we move more cautiously toward the new value (this can be useful if $\omega=1$ was unstable or to slow down oscillations in the iteration). Note that $\omega$ outside the range $(0,2)$ is generally not useful: $\omega \ge 2$ usually causes divergence, and $\omega = 0$ is trivial (no update at all).

% In matrix terms, the SOR iteration matrix can be written as $B_{\omega} = (D+\omega L)^{-1}[(1-\omega)D - \omega U]$ (which reduces to $(D+L)^{-1}U$ for $\omega=1$, consistent with Gauss–Seidel) . The freedom to choose $\omega$ allows one to tune the spectral radius $\rho(B_{\omega})$. In fact, SOR includes Jacobi as an extreme case: if we let $\omega \to 1$ from below, we approach Gauss–Seidel; if we consider $\omega$ very large, the effect is to almost ignore $L$ (since $D+\omega L \approx \omega L$ for large $\omega$) and one can show the iteration approaches the Jacobi method in the limit $\omega\to\infty$ . Thus $\omega$ “interpolates” between different schemes, and for each problem there is an optimal $\omega$ (within $(0,2)$) that minimizes $\rho(B_{\omega})$.

% Convergence and properties: Just as with Gauss–Seidel, SOR converges under broad conditions. If $A$ is symmetric positive definite, SOR converges for any $0 < \omega < 2$  (this classic result was proved by Ostrowski in 1947). In fact, $\omega=1$ (Gauss–Seidel) and $\omega=2$ are the endpoints of convergence for SPD matrices – the spectral radius $\rho(B_{\omega})$ as a function of $\omega$ is < 1 for $\omega\in(0,2)$ and hits 1 at $\omega=2$ . For strictly diagonally dominant matrices, SOR is also convergent for a wide range around $\omega=1$ (typically $(0,2)$ as well) . However, the optimal choice of $\omega$ can greatly affect the speed of convergence (it “stretches” or contracts the spectrum of $B_{\omega}$). If $\omega$ is too low, we are under-relaxing and convergence is safe but slower than necessary; if $\omega$ is too high (close to 2), we over-relax too much and the method may become unstable or sluggish.

% The challenge is that the optimal $\omega$ depends on the spectral properties of $A$ (or of the related iteration matrices). In some special cases, one can estimate it. For instance, if we denote $\mu = \rho(B_J)$ the spectral radius of the Jacobi iteration for $A$, and assume the eigenvalues of $B_J$ are real, theory shows the optimal $\omega$ satisfies roughly $\omega_{\text{opt}} \approx \frac{2}{1+\sqrt{1-\mu^2}}$ . The corresponding minimal spectral radius is $\rho(B_{\omega_{\text{opt}}}) = \frac{1-\sqrt{1-\mu^2}}{1+\sqrt{1-\mu^2}}$, which for small $\mu$ is about $\mu^2/4$ . In particular, when $\omega=1$ (Gauss–Seidel), one can show $\rho(B_{GS}) = \mu^2$ under these assumptions . With $\omega_{\text{opt}}$, the spectral radius is roughly $(\mu^2)/4$, making SOR about four times more efficient than Gauss–Seidel in asymptotic error reduction . In less technical terms: SOR can dramatically speed up convergence if $\omega$ is well-chosen. In practice, one might try a few values of $\omega$ to see which gives the fastest decrease in residual or error. Values $\omega>1$ (over-relaxation) typically give faster convergence up to a point, while $\omega<1$ might be used if the method is diverging or wildly oscillatory (to damp changes). It’s worth noting that SOR, like Gauss–Seidel, is sequential and dependent on ordering. Also, for $\omega \neq 1$ the iteration matrix $B_{\omega}$ is not symmetric even if $A$ is symmetric, which means SOR is not directly usable as a preconditioner for the Conjugate Gradient method (which requires symmetric positive definite preconditioners). In such cases, one often uses a Symmetric SOR (SSOR) scheme – essentially a forward SOR sweep followed by a backward SOR sweep in each iteration – which yields a symmetric positive definite effective preconditioner. (SSOR as a standalone iterative method doesn’t improve convergence much over SOR , but it is very useful as a preconditioner for nonstationary methods like CG.)

% Convergence Analysis: Spectral Radius and Norm Criteria

% As derived earlier, the convergence of any stationary method $x^{(k+1)} = B x^{(k)} + \mathbf{c}$ hinges on the spectral radius $\rho(B)$. If $\rho(B)<1$, the iteration converges for any initial guess, and if $\rho(B)\ge 1$, in general it will diverge or at best stagnate . The error after $k$ steps is $\boldsymbol{\epsilon}^{(k)} = B^k \boldsymbol{\epsilon}^{(0)}$, so $|\boldsymbol{\epsilon}^{(k)}| \le |B^k|\cdot|\boldsymbol{\epsilon}^{(0)}|$. Submultiplicativity gives $|B^k| \le |B|^k$, so a sufficient condition for convergence is $|B|<1$. In practice, one uses an induced matrix norm that reflects some known matrix structure. For example, if $A$ is strictly diagonally dominant by rows, one can use the infinity-norm: each row $i$ of $B_J = D^{-1}(L+U)$ has sum of absolute values $\frac{1}{|a_{ii}|}\sum_{j\neq i}|a_{ij}| < 1$, thus $|B_J|\infty<1$ which guarantees $\rho(B_J)\le |B_J|\infty < 1$ and convergence . Similarly, diagonal dominance by columns would ensure $|B|_1<1$. These norm criteria are often how convergence conditions are proved.

% For Jacobi and Gauss–Seidel methods, diagonal dominance of $A$ is a convenient sufficient condition for convergence (and for Gauss–Seidel, even a bit weaker condition can suffice, e.g. $A$ being an $H$-matrix with positive diagonal entries). In addition, symmetric positive definiteness of $A$ guarantees convergence of Gauss–Seidel (and SOR) but not of Jacobi  . SPD matrices have a nice property for Gauss–Seidel: one can show that $\rho(B_{GS})<1$ by an energy argument (Gauss–Seidel can be seen as a damped gradient descent for an SPD system, making it convergent). In contrast, Jacobi on an SPD matrix may fail if the matrix is not sufficiently diagonally dominant (as we saw in the counterexample). In summary:
% 	•	Jacobi method: Converges if $A$ is strictly (or irreducibly) diagonally dominant . It may also converge under other conditions (e.g. $A$ and $2D-A$ both SPD ), but not for all SPD matrices. If Jacobi converges, it is usually the slowest of the three methods.
% 	•	Gauss–Seidel method: Converges if $A$ is strictly (or irreducibly) diagonally dominant (same condition as Jacobi) and in many cases when $A$ is symmetric positive definite . In fact, any SPD matrix yields convergence of GS . Gauss–Seidel often converges in cases where Jacobi does not. If Jacobi converges, Gauss–Seidel will too, and typically faster .
% 	•	SOR method ($\omega$ in (0,2)): Converges for any symmetric positive definite $A$ if $0<\omega<2$ . Also converges for $A$ strictly diagonally dominant (with an appropriate range of $\omega$ which usually includes $(0,2)$) . In general, SOR shares convergence conditions with Gauss–Seidel when $\omega$ is in the safe range. However, an $\omega$ significantly greater than 1 can destabilize the iteration if $A$ doesn’t meet those strong conditions, so one must choose $\omega$ carefully. When $A$ is SPD, we have the guarantee that any $0<\omega<2$ works , and we can then worry about optimizing $\omega$ for speed.

% Finally, a note on rate of convergence: if $\rho(B)=\rho < 1$, the asymptotic convergence is linear and the error after $k$ iterations behaves like $|\boldsymbol{\epsilon}^{(k)}| \approx \rho^k |\boldsymbol{\epsilon}^{(0)}|$ for large $k$. The quantity $-1/\ln \rho$ is essentially the number of iterations required to reduce the error by a constant factor. A smaller spectral radius $\rho$ means faster convergence. One can also consider the spectral radius as the definitive measure: for example, if method A has $\rho_A = 0.5$ and method B has $\rho_B = 0.25$, then asymptotically method B will reduce the error roughly twice as fast as method A (since $0.25^k = (0.5^k)^2$). In the next section, we compare Jacobi, Gauss–Seidel, and SOR in these terms and discuss practical implications.



% Spectral radius vs. relaxation factor in SOR: The spectral radius $\rho(C_\omega)$ of the SOR iteration matrix is plotted as a function of the relaxation factor $\omega$, for different values of the Jacobi-method spectral radius $\mu = \rho(C_{\text{Jac}})$ (blue, red, and green curves correspond to different $\mu$) . SOR achieves its fastest convergence at the optimal $\omega$ (where each curve attains minimum $\rho$). For smaller Jacobi spectral radius (blue curve), the optimal $\omega$ yields a very small $\rho(C_\omega)$, illustrating that SOR can greatly accelerate convergence (spectral radius near 0) compared to Gauss–Seidel ($\omega=1$ point) or Jacobi (whose $\rho=\mu$ is much larger). Conversely, if Jacobi is barely convergent ($\mu$ close to 1, green curve), the gain from SOR is modest and excessive $\omega$ can be detrimental (note the rise of the green curve as $\omega \to 2$).

% The Jacobi, Gauss–Seidel, and SOR methods each have their pros and cons. In terms of convergence speed per iteration, Gauss–Seidel is almost always superior to Jacobi when both converge. In many cases (especially for SPD matrices), one can roughly estimate $\rho_{GS} \approx (\rho_{Jacobi})^2$ . For example, if $\rho_{Jacobi}=0.9$, then $\rho_{GS}\approx 0.81$; if $\rho_{Jacobi}=0.5$, then $\rho_{GS}\approx 0.25$. This squared improvement means Gauss–Seidel can cut the number of iterations dramatically. Empirically, if Jacobi takes $K$ iterations to reach a certain error tolerance, Gauss–Seidel might take on the order of $K/2$ iterations (or even better), though this is a rough rule of thumb. The SOR method can further improve the rate: with an optimally chosen $\omega$, the spectral radius of SOR can be much smaller. In the theoretical scenario mentioned earlier, $\rho_{SOR} \approx \rho_{GS}/4$ at optimal $\omega$ , meaning SOR might take on the order of $K/4$ iterations compared to Gauss–Seidel’s $K/2$ (hence an order of magnitude fewer iterations than Jacobi in that scenario ). In practice, the actual speedup depends on $A$ and how close to optimal your $\omega$ is. Even without exact tuning, SOR with $\omega=1.1$ or $1.2$ often outperforms plain Gauss–Seidel significantly if the matrix is well-behaved.

% However, practical considerations beyond iteration counts also matter:
% 	•	Work per iteration: Jacobi and Gauss–Seidel both require on the order of $\mathcal{O}(n^2)$ operations per iteration for a dense $n\times n$ system (or $\mathcal{O}(nz)$ for a sparse matrix with $nz$ nonzeros). In Gauss–Seidel and SOR, the work is organized as a forward substitution (triangular solve). In Jacobi, it’s a matrix-vector multiply with $L+U$. The cost is similar, though Gauss–Seidel’s sequential nature can limit parallel efficiency. Jacobi’s simultaneous updates can take advantage of parallel hardware (all $n$ updates independent), whereas Gauss–Seidel and SOR have inherent data dependencies . On modern computing systems, this means Jacobi can often achieve better parallel speed-up, despite needing more iterations, whereas Gauss–Seidel might not scale well with many processors.
% 	•	Memory usage: Jacobi needs to store two vectors (old and new $x$) if implemented in the straightforward way, since it cannot overwrite $x^{(k)}$ as it computes $x^{(k+1)}$ . Gauss–Seidel (and SOR) can update in place, requiring only one vector of storage for $x$. For very large problems, this memory saving might be relevant (though one extra vector is usually negligible compared to storing $A$).
% 	•	Convergence reliability: Jacobi is the simplest but also the most fragile – it may fail to converge even when the other methods do. Gauss–Seidel is more robust (converges on SPD matrices, for instance, where Jacobi might not). SOR introduces a parameter: a poorly chosen $\omega$ can slow convergence or even cause divergence if $\omega$ is outside the allowable range. In practice, one might start with $\omega=1$ (Gauss–Seidel) and then try increasing $\omega$ to speed up convergence, but with caution not to approach the instability threshold (for SPD matrices this threshold is $\omega=2$). If $\omega$ is slightly suboptimal, SOR will still converge (if it converges at all) but not as fast as possible. Thus, there is a tuning aspect: methods like Jacobi and Gauss–Seidel are parameter-free, whereas SOR demands choosing $\omega$. The optimal $\omega$ might be estimated if one has some spectral information, or found by trial.
% 	•	Ordering and flexibility: Gauss–Seidel and SOR depend on the ordering of unknowns. This can be a disadvantage if the “natural” ordering of the problem is not ideal for convergence. For example, interchanging equations can change $\rho(B_{GS})$. There are strategies like reordering the matrix to be closer to diagonal dominance or to reduce bandwidth, which can improve convergence. Jacobi’s convergence factor is invariant to reordering of equations (since it treats all equations independently in parallel), so it’s conceptually simpler in that regard. In SOR, the optimal $\omega$ can even depend on ordering. For certain structured matrices (like those arising from 2D Poisson’s equation on a grid), the classic result is that SOR with lexicographic Gauss–Seidel order and an optimal $\omega$ achieves the best possible linear convergence for a local method, but one can also use red-black ordering (a two-color ordering) to allow parallel Gauss–Seidel updates which is equivalent to a form of Jacobi or SOR on a checkerboard pattern.
% 	•	Usage as preconditioners: While Jacobi, Gauss–Seidel, and SOR can be used as solvers on their own, today they are often used as preconditioners or smoothers in more advanced iterative methods (like Conjugate Gradient or multigrid). For example, Gauss–Seidel (especially in its symmetric variant) is a common smoother in multigrid algorithms, and SSOR is a popular preconditioner for Conjugate Gradient on SPD systems . Jacobi (diagonal scaling) is sometimes used as a simple preconditioner as well. In these roles, the relative convergence speeds and ordering issues still matter, but one also considers factors like how well they reduce certain components of error (smoothing properties) and their cost per iteration. Gauss–Seidel has good smoothing for high-frequency error components in PDE problems, whereas Jacobi is cheaper per step on parallel hardware.

% In summary, Gauss–Seidel is generally superior to Jacobi in convergence per iteration (often roughly twice as fast in practice)  , at the cost of being sequential. SOR can be much faster than Gauss–Seidel if an effective relaxation factor $\omega$ is chosen, sometimes yielding an order of magnitude speed-up . However, SOR requires tuning and shares Gauss–Seidel’s sequential nature and sensitivity to ordering. All three methods are linear convergence (no “acceleration” in the sense of CG’s finite termination on SPD matrices), and for very large or ill-conditioned systems, even SOR may converge slowly, prompting the use of more sophisticated algorithms. Nevertheless, understanding Jacobi, Gauss–Seidel, and SOR is fundamental, as they illustrate how using partial information (diagonal or triangular part of $A$) can lead to iterative solvers, and they form building blocks for modern techniques. Their convergence conditions (diagonal dominance, SPD, etc.) are commonly checked to ensure the iterations will converge. In practice, one would choose based on the problem: Jacobi if parallelization is crucial and the matrix is well-conditioned or diagonally dominant; Gauss–Seidel if a bit more speed is needed and one can handle sequential updates; SOR if further speed is desired and some trial-and-error with $\omega$ is acceptable. And if even more speed is required, one might move beyond stationary methods to conjugate gradients or multigrid, possibly using these classical methods as components (preconditioners or smoothers) within those frameworks .





% \section{Iterative Non-Stationary Methods for Solving Linear Systems}

% Introduction

% Solving a linear system $A\mathbf{x}=\mathbf{b}$ is a fundamental task in computational mathematics. For small systems or well-behaved matrices, direct methods (like Gaussian elimination or LU factorization) can find exact solutions in finite steps. However, for large-scale or sparse systems (with millions of unknowns), direct solvers become prohibitively expensive in memory and time  . In such cases, one turns to iterative methods, which generate a sequence of approximations $\mathbf{x}^{(k)}$ converging to the true solution $\mathbf{x}^*$. These methods update an initial guess $\mathbf{x}^{(0)}$ repeatedly to reduce the residual $\mathbf{r}^{(k)} = \mathbf{b} - A\mathbf{x}^{(k)}$. Iterative algorithms are often the only feasible option for extremely large or complex systems arising from applications like discretized PDEs, optimization, or machine learning .

% Stationary vs. Non-Stationary Iterations: Iterative techniques fall into two broad categories: stationary methods and non-stationary methods  . In a stationary iterative method, the update formula is fixed and does not change with $k$. One can write the iteration as:
% \begin{equation}
    
% \mathbf{x}^{(k+1)} = B\,\mathbf{x}^{(k)} + \mathbf{c},
% \end{equation}

% where $B$ and $\mathbf{c}$ are fixed matrix and vector, independent of the iteration count . Examples include the Jacobi and Gauss–Seidel relaxations, where $B = M^{-1}N$ for some fixed splitting $A=M-N$ . These methods are simple and memory-efficient, but they often converge slowly and only under certain conditions (e.g. diagonal dominance or positive-definiteness of $A$)  . By contrast, non-stationary (or Krylov subspace) methods use information that changes at each iteration, such as the current residual, to adapt their next step . Their update formulas involve iteration-dependent coefficients (often obtained via inner products of residuals or search directions), making the process non-linear in the error. These methods are usually more sophisticated and effective, converging in far fewer steps than stationary schemes for difficult problems . For example, the conjugate gradient (CG) method – a non-stationary Krylov method for symmetric positive-definite matrices – can find the solution in a nearly optimal path, in stark contrast to simpler stationary approaches (like fixed-step steepest descent). This is illustrated in Figure 1, where the red path (CG) reaches the solution $\mathbf{x}$ in a direct zig-zag route, whereas the green path (gradient descent) makes slow orthogonal progress



% . In general, non-stationary methods harness more algebraic structure of $A$ at each step (such as previous search directions or residuals) to accelerate convergence, at the cost of a more complex analysis and implementation .

% Theoretical Foundation: Stationary vs. Non-Stationary Methods

% Stationary Iterative Methods. A linear stationary iteration for $A\mathbf{x}=\mathbf{b}$ can be derived by splitting $A = M - N$, with $M$ chosen for easy inversion . The iteration is then $M,\mathbf{x}^{(k+1)} = N,\mathbf{x}^{(k)} + \mathbf{b}$, or $\mathbf{x}^{(k+1)} = \mathbf{x}^{(k)} + M^{-1}(\mathbf{b} - A\mathbf{x}^{(k)})$ . Here, $M^{-1}N = I - M^{-1}A$ plays the role of a constant iteration matrix $C$ so that $\mathbf{e}^{(k+1)} = C,\mathbf{e}^{(k)}$ (where $\mathbf{e}^{(k)}=\mathbf{x}^{(k)}-\mathbf{x}^*$ is the error)  . Convergence is guaranteed if and only if all eigenvalues of $C$ lie strictly inside the unit circle, $\rho(C)<1$ . This spectral radius condition typically holds for matrices $A$ that are diagonally dominant, symmetric positive definite, or satisfy related criteria, after an appropriate splitting. Classic stationary methods include Jacobi’s method, Gauss–Seidel, and Successive Over-Relaxation (SOR)  . These rely on algebraic properties like matrix splittings and usually converge linearly (if at all), often requiring preconditioning or acceleration to be practical for large systems. Notably, stationary schemes have fixed update coefficients (e.g. a fixed relaxation factor $\omega$ in SOR), and do not “learn” from past iterations beyond the last iterate. This simplicity makes them easy to analyze and implement, but also limits their speed of convergence  .

% Non-Stationary Iterative Methods. Non-stationary methods abandon the fixed iteration matrix paradigm and instead adapt at each step. In these methods, the sequence of corrections is constructed based on dynamically computed coefficients, often involving inner products of current and previous residuals or search directions . Equivalently, one can view many non-stationary solvers as applying a varying higher-degree polynomial in $A$ to the initial residual at each iteration (instead of repeatedly applying a fixed $B$). The coefficients of this polynomial are chosen on the fly to optimize some criterion (such as residual norm). This adaptive approach allows non-stationary methods to accelerate convergence dramatically by canceling out error components associated with troublesome eigenvalues of $A$. Indeed, non-stationary Krylov methods can converge in at most $n$ steps in exact arithmetic by constructing a polynomial that annihilates the error (since the Krylov space eventually spans the whole $\mathbb{R}^n$)  . In practice, rounding errors and very large $n$ mean convergence in $\ll n$ steps is desired; the success of these methods hinges on the spectral properties of $A$. Intuitively, if $A$’s eigenvalues (or singular values) are clustered or well-conditioned, Krylov methods find a good approximation quickly by resolving large-scale components of the error first . Conversely, if $A$ has a wide spread of eigenvalues or pathological structure, convergence may stall or slow down (this is where preconditioning becomes crucial – discussed later).

% Algebraically, most non-stationary methods can be understood as projection methods onto a sequence of subspaces that grow with each iteration. Unlike stationary iterations that rely on the spectral radius of a fixed matrix $C$, Krylov methods leverage properties like orthogonality, optimality, and conjugacy in these subspaces. This makes their convergence analysis more involved (depending on polynomial approximations to $A^{-1}$ and the entire spectrum of $A$ ), but also more powerful. In summary, stationary methods are often used as preconditioners or smoothers (e.g. Gauss–Seidel in multigrid) rather than standalone solvers, whereas the heavy lifting for difficult linear systems is done by non-stationary Krylov solvers which adaptively exploit $A$’s algebraic structure.

% Krylov Subspace Methods: Orthogonality, Residual Minimization, and Projections

% Most non-stationary iterative methods of interest are Krylov subspace methods. The Krylov subspace of dimension $k$ generated by matrix $A$ and initial residual $\mathbf{r}^{(0)} = \mathbf{b}-A\mathbf{x}^{(0)}$ is defined as:

% \begin{equation}
    
% \mathcal{K}_k(A,\mathbf{r}^{(0)}) = \mathrm{span}\{\mathbf{r}^{(0)},\,A\mathbf{r}^{(0)},\,A^2\mathbf{r}^{(0)},\dots,A^{k-1}\mathbf{r}^{(0)}\}\,.
% \end{equation}

% This is the space of all linear combinations of the first $k$ powers of $A$ applied to the initial residual . Krylov methods iterate by expanding this subspace and seeking an improved solution $\mathbf{x}^{(k)}$ in the affine space $\mathbf{x}^{(0)} + \mathcal{K}_k(A,\mathbf{r}^{(0)})$. At each step, they enforce a condition that defines the new approximation, typically a projection condition: the residual at step $k$, $\mathbf{r}^{(k)} = \mathbf{b}-A\mathbf{x}^{(k)}$, is made orthogonal to some subspace (the test subspace). By construction $\mathbf{x}^{(k)}$ lies in the trial subspace $\mathbf{x}^{(0)} + \mathcal{K}_k$, so the method is choosing the “best” vector in that Krylov space according to some optimality criterion.

% Two common choices of projection lead to different families of methods: (1) Galerkin (Petrov-Galerkin) orthogonality: enforce $\mathbf{r}^{(k)}$ orthogonal to the trial subspace itself (or to a closely related subspace), and (2) Residual minimization: choose $\mathbf{x}^{(k)}$ that explicitly minimizes $|\mathbf{r}^{(k)}|$ over all candidates in $\mathbf{x}^{(0)} + \mathcal{K}_k$. The conjugate gradient (CG) method is an example of the first approach (Galerkin): for SPD matrices, CG chooses $\mathbf{x}^{(k)}$ such that the residual is orthogonal to $\mathcal{K}_k$ (which, in CG, also equals orthogonality to all previous residuals) . Equivalently, CG ensures the error $\mathbf{e}^{(k)}=\mathbf{x}^{(k)}-\mathbf{x}^*$ is minimized in the $A$-norm (an energy norm) within the trial subspace. On the other hand, GMRES (Generalized Minimal Residual) is the prototypical residual-minimization method: it selects $\mathbf{x}^{(k)}$ that minimizes $|\mathbf{r}^{(k)}|_2$ over $\mathbf{x}^{(0)}+ \mathcal{K}_k$ . This is achieved by building an orthonormal basis of $\mathcal{K}_k$ (using the Arnoldi process) and solving a small least-squares problem each iteration to obtain the coefficients that minimize the residual norm  . Other Krylov methods use variations or combinations of these two approaches – for example, BiCG uses a Petrov-Galerkin condition with two different subspaces, and BiCGStab adds a local residual-minimizing step to the BiCG process.

% Orthogonality and Conjugacy: Krylov solvers heavily exploit orthogonal or conjugate directions to improve efficiency. In CG, the successive residuals $\mathbf{r}^{(k)}$ are made orthogonal to each other (in the Euclidean inner product) and the search directions $\mathbf{p}^{(k)}$ are constructed to be $A$-conjugate (mutually orthogonal with respect to the $A$-inner product)  . This conjugacy condition ensures that each iteration corrects errors along a new eigen-direction of $A$ without spoiling the progress on previous directions. Thanks to these properties, CG implicitly builds an orthonormal basis of the Krylov subspace with short recurrences – only the last two residuals and search directions are needed to generate the next direction. In contrast, for a general non-symmetric $A$, one cannot maintain orthonormality with such short recurrences . GMRES addresses this by the Arnoldi process: it explicitly orthogonalizes the new Krylov vector $A\mathbf{v}^{(k)}$ against all previous basis vectors ${\mathbf{v}^{(1)},\dots,\mathbf{v}^{(k)}}$ to maintain an orthonormal basis of $\mathcal{K}_k$  . The benefit is a guarantee that the residual $\mathbf{r}^{(k)}$ is perpendicular to $\mathcal{K}_k$ (hence $ |\mathbf{r}^{(k)}|_2$ is minimized) . The cost, however, is that GMRES requires long recurrences: storing all $k$ basis vectors and performing $O(k)$ work to orthogonalize each new vector, which becomes expensive as $k$ grows  .

% Other Krylov methods find compromises. BiCG (Bi-Conjugate Gradient) generates two coupled sequences of basis vectors: one for the Krylov subspace of $A$ (using the residual $\mathbf{r}^{(k)}$) and one for the Krylov subspace of $A^T$ (using a shadow residual $\tilde{\mathbf{r}}^{(k)}$) . It enforces a bi-orthogonality: the residuals from the $A$-sequence are orthogonal to the $A^T$-Krylov basis (which can be seen as a Galerkin condition with a different test subspace). This yields a short-recurrence method for non-symmetric $A$—like CG, only a few vectors are kept—but at the expense of potential breakdown (if the process encounters a situation where an inner product vanishes unexpectedly) . Variants like QMR (Quasi-Minimal Residual) apply additional techniques (least-squares smoothing of BiCG’s updates) to avoid breakdowns and smooth the convergence  . CGS (Conjugate Gradient Squared), described below, takes a different approach: it essentially squares the BiCG polynomial, eliminating the need for the $A^T$ sequence by merging it with the $A$-sequence . And BiCGStab (BiCG Stabilized) further refines BiCG/CGS by adding a one-step residual minimization (a small GMRES step of length 1) at each iteration to stabilize convergence  .

% Residual Minimization and Optimality: A hallmark of Krylov methods is that many of them are optimal in some sense. GMRES is $l_2$-optimal: at iteration $k$ it produces the smallest possible $|\mathbf{r}^{(k)}|_2$ among all linear combinations of ${I, A, \dots, A^{k-1}}$ applied to $\mathbf{r}^{(0)}$ . In other words, GMRES finds the best approximation in the $k$-dimensional Krylov space by a least-squares criteria. MINRES is similar for symmetric systems (minimizing residual in the $A$-norm for indefinite symmetric $A$). Conjugate Gradient is optimal in the $A$-norm (it minimizes the error norm $|\mathbf{e}^{(k)}|_A$ over the Krylov space for SPD $A$). This optimality leads to predictable convergence bounds: for CG one can prove

% \begin{equation}
    
% \|\mathbf{e}^{(k)}\|_A \;\le\; 2\left(\frac{\sqrt{\kappa(A)}-1}{\sqrt{\kappa(A)}+1}\right)^k \,\|\mathbf{e}^{(0)}\|_A,
% \end{equation}

% where $\kappa(A)$ is the condition number of $A$ . This shows the error drops at least as fast as $(\sqrt{\kappa}-1)/(\sqrt{\kappa}+1)$ per iteration (CG often converges faster than this worst-case bound due to spectral clustering or superlinear convergence when eigenvalues cluster ). GMRES lacks such a simple bound in terms of $\kappa(A)$, because for non-normal matrices the convergence can depend on the entire spectrum distribution and eigenvector conditioning. Nonetheless, GMRES will converge in at most $n$ steps (exact in $\mathbb{R}^n$) , and in practice often much sooner if the polynomial of best approximation of $A^{-1}$ on the spectrum is of low degree. Non-stationary methods are typically superior to stationary ones because they build these optimal (or near-optimal) polynomials adaptively, whereas a stationary method effectively applies the same fixed polynomial (the iteration matrix $C$) repeatedly, without targeting the troublesome spectral components of the error.

% Key Examples of Non-Stationary Iterative Methods

% Conjugate Gradient (CG) – prototype for SPD systems

% CG is the oldest and one of the most powerful non-stationary methods, applicable to symmetric positive-definite (SPD) matrices. It generates a sequence of conjugate (A-orthogonal) directions $\mathbf{p}^{(k)}$ and ensures the residuals $\mathbf{r}^{(k)}$ are orthogonal to each other . Each iteration chooses step lengths $\alpha_k$ and $\beta_k$ (via inner products of residuals) to minimize the quadratic error (since solving $A\mathbf{x}=\mathbf{b}$ for SPD $A$ is equivalent to minimizing a convex quadratic function). As a result, CG converges in at most $n$ steps, and in practice often in far fewer: the error after $k$ steps satisfies the bound above, which depends on $\sqrt{\kappa(A)}$ . If $A$’s eigenvalues cluster, convergence is significantly accelerated (sometimes superlinear). CG’s algorithmic structure is simple: each iteration does one matrix-vector multiply $A\mathbf{p}^{(k)}$, two inner products (to compute $\alpha_k$ and update $\beta_{k}$), and a few vector saxpy operations. Storage is minimal – it needs only a handful of $n$-vectors (for $\mathbf{x}^{(k)}$, $\mathbf{r}^{(k)}$, $\mathbf{p}^{(k)}$, and $A\mathbf{p}^{(k)}$), making it very memory-efficient. This combination of speed and low memory makes CG extremely effective for large SPD problems . Indeed, CG is often the method of choice for symmetric linear systems arising from e.g. finite element discretizations of elliptic PDEs (structural mechanics, diffusion problems, etc.), especially when a good preconditioner is available. The theoretical elegance of CG (polynomial optimality and three-term recurrence) set the stage for all later Krylov methods.

% Generalized Minimal Residual (GMRES) – robust solver for unsymmetric systems

% GMRES (Saad and Schultz, 1986) is a widely used method for general non-symmetric (or non-Hermitian) matrices. GMRES constructs an orthonormal basis ${v_1,\dots,v_m}$ for $\mathcal{K}_m(A,r^{(0)})$ via the Arnoldi process , and computes the new approximation $\mathbf{x}^{(m)} = \mathbf{x}^{(0)} + \sum_{j=1}^m y_j,v_j$ that minimizes $|\mathbf{b}-A\mathbf{x}^{(m)}|_2$. This leads to solving a small $m\times m$ least-squares problem (often written as $H \mathbf{y} \approx \beta e_1$ for the upper Hessenberg matrix $H$ obtained from Arnoldi)  . By design, GMRES produces the smallest residual norm possible after $m$ steps (among all Krylov methods of that subspace dimension) . This often gives it excellent robustness – GMRES will converge whenever the underlying Krylov subspace contains a good approximation, and it is guaranteed not to stagnate before reaching the true solution (except in pathological cases where the best polynomial approximation of $A^{-1}$ is the trivial one of degree $n$). In exact arithmetic GMRES would solve the system in $\le n$ steps ; in practice one uses restarted GMRES, denoted GMRES($m$), which limits the basis size to $m$ to control storage and orthogonalization cost  . For example, GMRES(30) restarts every 30 iterations (discarding the accumulated basis and starting anew from the current residual). Restarts prevent runaway memory use, but if chosen too small can hurt convergence or cause stagnation (failure to improve until $m$ is large enough)  . GMRES without restart has $O(nm + m^2)$ work per $m$ iterations (due to orthogonalization) and $O(nm)$ storage, which is why $m$ is kept moderate in practice.

% Advantages: GMRES is a “black-box” solver – it works on any square matrix (even singular ones in certain cases, with appropriate handling) and is often very stable. It excels for tough nonsymmetric problems where other methods (BiCG or CGS) might break down or exhibit wild residual oscillations. Disadvantages: The memory and CPU cost of the growing Arnoldi process. For very large $n$, GMRES($m$) with small $m$ might converge slowly, whereas large $m$ is infeasible. Strategies like adaptive restarting or augmenting the subspace with known troublesome vectors (e.g. eigenvector approximations) are used to mitigate this. Despite these issues, GMRES remains a workhorse in industrial solvers for problems like computational fluid dynamics (Navier–Stokes equations), where matrices are large, sparse, and non-symmetric, and a direct solve is impossible.

% Bi-Conjugate Gradient (BiCG) – extending CG to non-symmetric systems

% BiCG is a Krylov method that generalizes the conjugate gradient concept to matrices $A$ that are not symmetric. Since CG’s simplicity relies on the symmetry of $A$ (which guarantees a single Krylov sequence suffices), BiCG tackles the general case by generating two sequences of vectors: ${p_k}$ in the column space of $A$ and ${\tilde{p}_k}$ in the column space of $A^T$. It simultaneously advances these sequences such that they remain bi-orthogonal (i.e. $\tilde{p}_i^T (A,p_j) = 0$ for $i\neq j$), ensuring a Petrov-Galerkin condition: the residual $r^{(k)}$ is made orthogonal to the space spanned by ${\tilde{p}_1,\dots,\tilde{p}_k}$ . In practice, one can implement BiCG with a short two-term recurrence similar to CG, but requiring two matrix-vector multiplications per iteration: one with $A$ and one with $A^T$ . The need to multiply by $A^T$ is a drawback if $A$ is not readily transposable (e.g. when $A$ is given as a black-box operator or in a distributed format). BiCG’s cost per iteration is roughly twice that of CG (two matvecs, and a few dot-products) and it uses a small, fixed number of vectors for storage (on the order of 4–6 vectors of length $n$). Like CG, BiCG does not explicitly minimize any residual norm; thus its convergence can be irregular – residuals might fluctuate in norm. When it converges, BiCG often achieves a solution of accuracy comparable to GMRES in fewer iterations (but with each iteration costing more) . However, BiCG can suffer breakdown if at some step the direction vectors lose bi-orthogonality (denominator in recurrence becomes zero). Look-ahead strategies exist to overcome breakdowns at additional cost , but other methods were developed to improve upon BiCG’s reliability.

% Despite its pitfalls, BiCG is historically important and forms the backbone for more refined solvers. For instance, on SPD matrices BiCG reduces to CG (and indeed delivers identical results as CG, but at twice the work, making it inefficient for SPD cases) . BiCG is seldom used on its own today; instead one of its variants is preferred, as discussed next.

% Conjugate Gradient Squared (CGS) – squaring the BiCG polynomial for speed

% CGS (Sonneveld, 1989) was introduced to avoid the explicit computation with $A^T$ in BiCG and potentially double the convergence speed. The idea is clever: in BiCG, one can view the residual after $k$ steps as $r^{(k)} = P_k(A),r^{(0)}$ for some $k$-degree polynomial $P_k$. CGS considers $P_k^2(A),r^{(0)} = P_k(A),P_k(A),r^{(0)}$, effectively using the same polynomial for both $A$ and (implicitly) $A^T$ sequences . By interleaving the BiCG recurrences, CGS produces a new residual sequence that (in exact arithmetic) corresponds to the squared minimal polynomial of BiCG, ideally achieving the solution in roughly half as many iterations. In practice, CGS often does converge in fewer iterations than BiCG (roughly, it may finish in $k/2$ iterations compared to BiCG’s $k$) . Importantly, CGS does not require $A^T$ – it uses only $A$-multiplications, performing two of them per iteration (so the flops per iteration are similar to BiCG’s two matvecs, but without needing the transpose)  . This made CGS attractive for problems where forming $A^T$ is inconvenient.

% The major drawback of CGS is that squaring the polynomial can lead to highly irregular convergence . The residual norm in CGS sometimes oscillates wildly, even increasing temporarily or failing to decrease monotonically. This can cause loss of accuracy or instability in finite precision. In some cases, CGS can stagnate or diverge where BiCG would have converged steadily. Nonetheless, CGS was an important development because it eliminated the $A^T$ requirement and spurred research into fixes for the erratic convergence. Today, CGS is generally supplanted by BiCGStab (or variants like BiCGStab(l)), which we discuss next, but its idea of polynomial squaring is conceptually valuable.

% Bi-Conjugate Gradient Stabilized (BiCGStab) – smoothing and improving BiCG/CGS

% BiCGStab (van der Vorst, 1992) is one of the most popular general-purpose solvers for unsymmetric systems. It was designed specifically to cure the weaknesses of CGS while retaining its strengths  . BiCGStab starts from the BiCG two-sequence framework (like CGS does) but introduces an additional “stabilization” step in each iteration. In each BiCGStab iteration, one first performs a BiCG-like update (one $A$ matvec, producing an intermediate residual as in BiCG), and then follows up with a single-step GMRES minimization (one more $A$ matvec) restricted to the space spanned by the current residual and the $A$-applied residual . In effect, BiCGStab applies a degree-1 polynomial smoothing $Q_i(A)$ to the result of the BiCG iteration $P_i(A)$ . The smoothing polynomial $Q_i$ is chosen (via an inner product) to minimize the residual norm in that 2-dimensional subspace (often denoted by an $\omega_i$ parameter in the algorithm). The outcome is that residuals are much smoother – BiCGStab tends to avoid the extreme oscillations of CGS . In many cases BiCGStab’s convergence is also more reliable than BiCG, since it combines the fast convergence modes of BiCG with occasional damping.

% Each iteration of BiCGStab costs two $A$–multiplications (like CGS and BiCG) and a handful of dot-products and axpy operations. It also does not require $A^T$ (like CGS, it only uses $A$ on-the-fly). The storage overhead is small (a few vectors, similar to BiCG). Thus, per iteration cost and memory of BiCGStab are comparable to BiCG/CGS . The benefit is that BiCGStab often converges as fast or faster than CGS in practice and with far fewer erratic spikes in residual norm  . It is regarded as one of the best general-purpose Krylov solvers for nonsymmetric matrices when one wants a good balance of convergence speed, robustness, and low memory. There are further improvements like BiCGStab($\ell$) which apply an $\ell$-degree GMRES smoothing polynomial instead of 1 (giving more robust but sometimes slower convergence), and IDR(s) (Induced Dimension Reduction) methods that have fixed subspace recycling, but BiCGStab remains highly competitive for many problems.

% Summary of Algorithmic Structure: The table below summarizes and compares some features of these key methods:
% 	•	CG: short recurrences (2-term), one $A$ matvec per iteration; optimal for SPD matrices; storage ~4 vectors; convergence rate tied to $\sqrt{\kappa(A)}$.
% 	•	GMRES: long recurrences (full orthogonalization), one $A$ matvec per iteration; optimal residual in 2-norm; storage grows with iterations (or with restart length $m$); robust for general $A$ but cost grows and may need restarting.
% 	•	BiCG: short recurrences (2-term) but two coupled sequences; two matvecs ($A$ and $A^T$) per iteration; not norm-minimizing, residual can oscillate; storage ~6 vectors; works for general $A$ but can break down.
% 	•	CGS: short recurrences, no $A^T$ needed; two $A$ matvecs per iteration (like BiCG); often faster convergence in iterations than BiCG but with highly irregular residual behavior; storage ~6 vectors; rarely used now due to instability.
% 	•	BiCGStab: short recurrences, no $A^T$; effectively two $A$ matvecs per iteration; smoother convergence than CGS (and often BiCG); storage ~6–8 vectors; one of the most reliable choices for general unsymmetric systems.

% All these methods (except GMRES without restart) have $O(n)$ storage and $O(n)$ work per iteration plus the cost of matvecs (which is $O(\text{nnz}(A))$ for sparse $A$, or more if $A$ is dense). In modern high-performance computing, the matrix-vector multiply is usually the dominant cost, so methods with more matvecs per iteration (BiCG, CGS, BiCGStab) are a bit more expensive per step than ones with one matvec (CG, GMRES). However, they may converge in fewer iterations. There is also a consideration of communication overhead: methods like CG and BiCGStab require a couple of global inner products per iteration, which can limit scalability on parallel machines. Research into s-step Krylov methods and avoiding global synchronization is ongoing, but beyond our scope here.

% Convergence Properties and the Role of Preconditioners

% The convergence of non-stationary iterative methods is intricately linked to the spectrum of $A$ (its eigenvalues, and for non-normal $A$, the behavior of the pseudoeigenvalues as well). Roughly speaking, Krylov methods struggle when $A$ is ill-conditioned or has eigenvalues spread over a large range or lying on problematic locations in the complex plane. For example, CG’s error bound involves $\kappa(A) = \lambda_{\max}/\lambda_{\min}$; if this condition number is large, convergence is slow . GMRES’s performance can degrade if $A$ has eigenvalues on or near the unit circle or if $A$ is nearly defective (non-diagonalizable), causing the Krylov basis vectors to become almost linearly dependent. One often observes stagnation or extremely slow improvement in those cases. Convergence analysis is generally hard – for CG there is a rich theory involving Chebyshev polynomials and A-norm energy minimization, for GMRES and BiCG there are results relating residual polynomials to Faber or Leja points in the complex plane , but precise convergence rates are usually problem-specific.

% A few general statements can be made:
% 	•	Krylov methods will converge in $n$ steps in exact arithmetic (since $\deg p(A) \le n$ for the minimal polynomial)  , but round-off error effectively limits the useful subspace dimension and can cause stagnation before $n$ steps. Typically, one desires convergence in $k \ll n$ steps.
% 	•	Eigenvalue distribution: If the eigenvalues of $A$ cluster in a few intervals, a low-degree polynomial can approximate $1/A$ well on that spectrum, leading to fast convergence. Conversely, if eigenvalues are widespread, the residual polynomial needs higher degree. For CG, tightly clustered eigenvalues (except a few outliers) lead to superlinear convergence as the method “zeros out” whole clusters quickly. For GMRES/BiCG, a helpful metric is the field of values (numerical range) – if that is narrow or well-separated from the origin, convergence is faster. Indefinite or sign-changing spectra (for which some eigenvalues are very small or zero) pose challenges – CG will stagnate if any $\lambda_i=0$ (singular $A$) or if $\kappa$ is huge, etc.
% 	•	Non-normal matrices: Even if the eigenvalues suggest easy convergence, a highly non-normal matrix can cause slow convergence (GMRES in particular) because eigenvectors are not orthogonal and the Krylov basis may need to be large to capture the difficult directions. In such cases, residual norms don’t tell the full story of error (someone might see “delay” in convergence until certain invariant subspaces are approximated).

% To address slow convergence, one almost always employs a preconditioner $M$. A preconditioner is an auxiliary matrix (or operator) that approximates $A^{-1}$ in some sense, used to transform the linear system into one that is more favorable spectrally . Commonly $M$ is chosen such that $M^{-1}A$ (or $AM^{-1}$) has a smaller condition number or more clustered eigenvalues than $A$. Krylov methods can then be applied to the preconditioned system (e.g. solve $M^{-1}A x = M^{-1}b$ by iterations). Good preconditioners drastically improve convergence – for instance, a well-tuned Incomplete LU (ILU) or multigrid preconditioner can reduce the needed iterations by an order of magnitude or more. In effect, preconditioning incorporates some of the benefits of stationary methods (which split $A$ = preconditioner + remainder) into Krylov solvers . Indeed, one can view many Krylov methods as accelerated stationary iterations: e.g. preconditioned CG is like a polynomially accelerated form of the Jacobi or Gauss–Seidel method used as $M^{-1}$ . Preconditioning is a vast topic on its own – key points are that the preconditioner must be computationally cheap to apply (ideally one solve of $M$ per iteration, costing much less than solving $A$ directly) and should make $M^{-1}A$ spectrally friendly (clustered eigenvalues, low $\kappa(M^{-1}A)$, or deflated problematic modes).

% Convergence in practice: With a decent preconditioner, CG typically converges in $\mathcal{O}(\sqrt{\kappa(M^{-1}A)})$ iterations, often a small multiple of the problem’s dimension in the context of PDEs (e.g. for elliptic PDEs, multigrid can make iterations independent of $n$). GMRES and BiCGStab with preconditioning often show convergence profiles where the residual norm drops steadily (sometimes with a few stagnation plateaus that eventually resolve). One must be cautious: methods like BiCG and BiCGStab can still break down (e.g. if $M^{-1}A$ has certain eigenvalue distributions causing the BiCG formulas to hit a zero divisor). Variants like QMR or GMRES-based methods are used if that is a concern. In summary, convergence is assured for SPD problems (CG always converges or finds an exact solution in $n$ steps) and generally expected for reasonably well-conditioned non-symmetric problems. For very difficult matrices (highly ill-conditioned or nearly singular), iterative solvers might converge too slowly or fail, and one might need to resort to direct solvers or specialized techniques.

% Algorithmic Considerations: Updates, Storage, and Cost

% We have touched on many algorithmic aspects of these methods, but we summarize the key points:
% 	•	Update rules & recurrences: Non-stationary methods typically use recurrence formulas to update the solution and residual without storing the entire Krylov basis (except GMRES). For CG:
% $$\alpha_k = \frac{\mathbf{r}^{(k-1)T}\mathbf{r}^{(k-1)}}{\mathbf{p}^{(k)T}A,\mathbf{p}^{(k)}}, \qquad \mathbf{x}^{(k)} = \mathbf{x}^{(k-1)} + \alpha_k \mathbf{p}^{(k)}, \qquad \mathbf{r}^{(k)} = \mathbf{r}^{(k-1)} - \alpha_k A\mathbf{p}^{(k)}$$
% then $$\beta_k = \frac{\mathbf{r}^{(k)T}\mathbf{r}^{(k)}}{\mathbf{r}^{(k-1)T}\mathbf{r}^{(k-1)}}, \qquad \mathbf{p}^{(k+1)} = \mathbf{r}^{(k)} + \beta_k \mathbf{p}^{(k)}.$$
% This 2-term recurrence keeps CG efficient  . BiCG has a similar pair of recurrences for $\mathbf{p},\tilde{\mathbf{p}}$ sequences, involving inner products between $r$ and $\tilde{r}$ for $\beta_k$ updates. BiCGStab’s update is more complex (two-phase: first like BiCG to get an interim $r’$ and then an $\omega$ step for minimization), but it still uses only a fixed number of vectors. GMRES, by contrast, updates via the Arnoldi relation $A V_m = V_{m+1} H_m$ (where $V_m$ has orthonormal basis vectors as columns, and $H_m$ is Hessenberg). The solution update in GMRES comes from solving $H_m \mathbf{y} = \beta e_1$ and setting $\mathbf{x}^{(m)} = \mathbf{x}^{(0)} + V_m \mathbf{y}$. After $m$ steps, if not converged, GMRES may restart by discarding $V_m$ except perhaps the last residual vector.
% 	•	Storage requirements: As noted, CG / BiCG / CGS / BiCGStab require storing only a fixed number of vectors (on the order of 5–8 vectors of length $n$). For example, BiCGStab needs $\mathbf{r}, \mathbf{r}_0$ (initial residual copy for BiCG orthogonality), $\mathbf{p}, \tilde{\mathbf{p}}$ (search directions), and a few temporary vectors like $A\mathbf{p}$ and $A\tilde{\mathbf{p}}$ or intermediate $u$ vectors – in total maybe 6–7 $n$-length arrays. This is very modest in modern memory terms, allowing solutions of extremely large systems (millions of unknowns) as long as one can perform the matrix-vector product in distributed memory. GMRES(m) needs to store $m+1$ basis vectors plus the $m\times m$ Hessenberg (and some smaller work vectors). If $m=30$, this is 31 vectors – which for $n$ large could be significant but still feasible; if one tried GMRES with no restart on a million-unknown problem for 1000 iterations, storing 1000 vectors of length $10^6$ would be infeasible (that’s $10^9$ entries). Thus, storage is a critical factor in choosing GMRES vs. a BiCG-type method. Methods like MINRES, QMR also use short recurrences (MINRES stores 3 vectors, QMR a few more but still $O(1)$).
% 	•	Computational cost per iteration: All these methods are dominated by the cost of sparse matrix-vector multiplications. If $A$ has $\nnz$ nonzeros, one matvec is $O(\nnz)$ operations. CG does 1 matvec per iteration, BiCG/CGS/BiCGStab do 2 matvecs, GMRES does 1 matvec but also $O(mn)$ orthonormalization work each step. In addition, each method performs a few vector dot products (inner products) and SAXPYs (vector additions) per iteration. Typically: CG has 2 inner products and 3 saxpy; BiCGStab has ~4 inner products (including those for BiCG and for the $\omega$ computation) and ~6 saxpy; GMRES has 1 inner product per orthogonalization step plus solving the small least-squares (negligible) and one or two inner products to check convergence, etc. In a single-core model, these vector operations are low-order compared to matvec, but in parallel, global inner products are costly due to communication synchronization. One notable difference is that BiCG and CG allow some overlap of computation (e.g. one can sometimes overlap the second matvec with communication of inner product results), whereas GMRES’s Arnoldi process is more sequential (one cannot start the next matvec until orthonormalization of the current vector is done). This means in practice, on parallel clusters, CG and BiCGStab may scale better (less communication bottleneck) than GMRES for very large runs. There are communication-avoiding GMRES variants that address this.
% 	•	Breakdowns and numerical stability: Non-stationary methods must handle the possibility of breakdown. CG can break down only if $A$ is not SPD (or if hit by floating-point issues like divide-by-zero due to $\mathbf{r}^{(k)}$ becoming nearly 0 unexpectedly – which actually indicates convergence). BiCG and derivatives can suffer pivot breakdown (if $\tilde{r}^{(0)T}r^{(k)} = 0$ for some $k$ before convergence) or lucky breakdown (converged early). QMR avoids breakdown by lookahead; GMRES doesn’t break down (unless one hits the exact solution and residual becomes zero, which is a “happy breakdown” meaning it solved early). Numerical stability: all these methods are backward stable in a mild sense, but long recurrence methods like GMRES can accumulate orthogonalization error – one may need reorthogonalization to maintain the basis quality. Short recurrence methods like CG don’t have this issue (the three-term recurrence is stable for SPD $A$), but BiCG and CGS can be less numerically stable (because they don’t explicitly orthogonalize; they rely on ideal orthogonality that can drift in floating point). BiCGStab tends to be more stable than BiCG/CGS in practice but can still suffer if $A$ is very ill-conditioned (one might then restart or switch to GMRES).

% Applications and Use Cases

% Non-stationary iterative solvers are preferred in virtually all large-scale applications where direct methods fail or are too slow. Some prominent scenarios:
% 	•	Scientific Simulations (PDEs): Solving linear systems arising from finite difference/element/volume discretizations of PDEs is a prime domain. For example, in structural analysis or thermal diffusion (with SPD stiffness or conductivity matrices), CG with a multigrid or ILU preconditioner is the gold standard – it can solve systems with millions of unknowns where direct factorization is infeasible. In fluid dynamics (Navier–Stokes equations, which give non-symmetric matrices due to convection terms or saddle-point systems), GMRES or BiCGStab are commonly used, often with domain-decomposition or multigrid preconditioners  . These methods allow engineers to simulate aerodynamic flows, weather models, or groundwater movement by solving large linear systems at each time step or nonlinear iteration.
% 	•	Computer Graphics and Vision: Large linear systems appear in solving for instance image deblurring (where matrices might be ill-conditioned blurring operators), or in graphics for solving Poisson equations (again CG or MINRES are used, often on GPUs for real-time). Krylov methods with FFT-based preconditioners can handle huge image processing tasks efficiently, where naive methods would be too slow.
% 	•	Optimization and Machine Learning: Conjugate gradient is also an algorithm for quadratic optimization; in machine learning, CG can be used to solve normal equations or as a subroutine in Newton’s method for large logistic regression, etc. While first-order methods dominate ML, CG is sometimes used for its fast convergence property on quadratic subproblems. The Hessians in such problems are SPD, making CG applicable. Nonlinear conjugate gradients (an extension of CG) are widely used in large-scale optimization because they avoid storing a full Hessian but still accelerate convergence over basic gradient descent.
% 	•	Iterative Refinement and Precision Engineering: Sometimes a direct method is used to get an approximate solution and then an iterative method (like GMRES) refines the solution to high precision by solving the correction equation. GMRES and BiCGStab are useful here because they cope with the general nonsymmetric systems that can arise in, say, least-squares problems or other correction equations.
% 	•	Parallel Computing: Importantly, non-stationary methods are much more amenable to parallelization than stationary Gauss–Seidel or SOR. Stationary methods like Gauss–Seidel are inherently sequential (each step uses the latest components of $\mathbf{x}$ as soon as they are available), which complicates parallel implementation. Krylov methods use mostly matrix-vector multiplications and vector updates, which are highly parallelizable across distributed memory or GPUs. As a result, modern supercomputing simulations exclusively rely on Krylov solvers (with appropriate preconditioners) to utilize thousands of processors simultaneously  . For instance, a petascale simulation might solve a sparse linear system with $10^9$ unknowns using a preconditioned CG or GMRES solver running on thousands of cores – something utterly impossible without iterative methods.
% 	•	When are stationary methods still used? Typically as smoothers in multigrid or as simple preconditioners. For example, one might use a few Gauss–Seidel sweeps as a preconditioner inside CG to damp high-frequency error components (a common multigrid practice). But for obtaining a high-accuracy solution, non-stationary methods are usually preferred due to their superior convergence properties.

% In conclusion, iterative non-stationary methods like CG, GMRES, BiCG, BiCGStab, and others form a rich mathematical framework for solving linear systems by projection onto Krylov subspaces. They leverage algebraic properties (orthogonality, conjugacy, residual minimization) to achieve fast convergence, especially when aided by good preconditioners. Each method has trade-offs in memory, computation per iteration, and applicability (depending on $A$ being symmetric or not, etc.), but together they enable the solution of immense linear systems that would be intractable by direct methods. Ongoing research continues to improve these algorithms – for example, by making them more resilient to rounding errors, reducing communication, or extending them to block and multi-right-hand-side cases – ensuring that Krylov solvers remain at the forefront of numerical linear algebra for years to come.

% Sources: Key references for this overview include Barrett et al.’s Templates for the Solution of Linear Systems  , the classic book by Y. Saad Iterative Methods for Sparse Linear Systems  , and various lecture notes and Wikipedia articles on Krylov subspace methods   which provide detailed algorithmic descriptions and theoretical results.